\hypertarget{python-3.8-to-3.10-type-hinting-protocols-and-structural-subtyping}{%
\chapter{Python 3.8 to 3.10: Type Hinting Protocols and Structural
Subtyping}\label{python-3.8-to-3.10-type-hinting-protocols-and-structural-subtyping}}

\hypertarget{nominal-vs.-structural-subtyping}{%
\subsection{16.1 Nominal vs.~Structural
Subtyping}\label{nominal-vs.-structural-subtyping}}

\begin{itemize}
\tightlist
\item
  \textbf{Nominal Subtyping}: Standard object-oriented inheritance. An
  class \passthrough{\lstinline!A!} is a subclass of
  \passthrough{\lstinline!B!} only if \passthrough{\lstinline!A!}
  explicitly inherits from \passthrough{\lstinline!B!}
  (\passthrough{\lstinline!class A(B):!}).
\item
  \textbf{Structural Subtyping (Protocols / PEP 544)}: Also known as
  static duck typing. A class \passthrough{\lstinline!A!} is compatible
  with a Protocol \passthrough{\lstinline!P!} if
  \passthrough{\lstinline!A!} implements all methods and variables
  defined in \passthrough{\lstinline!P!} with matching signatures,
  regardless of explicit inheritance declarations.
\end{itemize}

\begin{center}\rule{0.5\linewidth}{0.5pt}\end{center}

\hypertarget{runtime-protocols-runtime_checkable}{%
\subsection{\texorpdfstring{16.2 Runtime Protocols \&
\texttt{@runtime\_checkable}}{16.2 Runtime Protocols \& @runtime\_checkable}}\label{runtime-protocols-runtime_checkable}}

By default, PEP 544 protocols are strictly compile-time constructs
erased during execution. However, annotating a protocol with the
\passthrough{\lstinline!@runtime\_checkable!} decorator permits the use
of \passthrough{\lstinline!isinstance()!} and
\passthrough{\lstinline!issubclass()!} checks at runtime.

\hypertarget{internal-hook-overrides}{%
\subsubsection{1. Internal Hook
Overrides}\label{internal-hook-overrides}}

The \passthrough{\lstinline!@runtime\_checkable!} decorator replaces the
standard class methods for instance validation. It overrides the
metaclass dunder methods: *
\passthrough{\lstinline!\_\_instancecheck\_\_!}: Called when performing
\passthrough{\lstinline!isinstance(obj, ProtocolClass)!}. *
\passthrough{\lstinline!\_\_subclasscheck\_\_!}: Called when performing
\passthrough{\lstinline!issubclass(SubClass, ProtocolClass)!}.

\hypertarget{runtime-evaluation-loop}{%
\subsubsection{2. Runtime Evaluation
Loop}\label{runtime-evaluation-loop}}

The CPython implementation of
\passthrough{\lstinline!\_\_subclasscheck\_\_!} on a runtime checkable
protocol executes the following logic:

\begin{lstlisting}[language=Python]
# Conceptual runtime protocol verification
def __subclasscheck__(cls, subclass):
    if not isinstance(subclass, type):
        raise TypeError("issubclass() arg 1 must be a class")
    
    # Iterate over all defined attributes in the protocol
    for attr in cls._protocol_attrs:
        # Check if the subclass or its MRO implements the attribute
        if not any(attr in s.__dict__ for s in subclass.__mro__):
            return False
    return True
\end{lstlisting}

\begin{quote}
{[}!WARNING{]} Runtime protocol checks only verify the
\textbf{existence} of attributes, not their method signatures, types, or
parameters. Furthermore, traversing the entire MRO of a subclass for
every attribute check incurs significant runtime overhead compared to
simple nominal class validations.
\end{quote}

\begin{center}\rule{0.5\linewidth}{0.5pt}\end{center}

\hypertarget{static-type-checking-and-variance}{%
\subsection{16.3 Static Type Checking and
Variance}\label{static-type-checking-and-variance}}

Static type checkers (such as Mypy or Pyright) compile type annotations
into a directed graph of type variables and interfaces.

\hypertarget{mathematical-definitions-of-variance}{%
\subsubsection{1. Mathematical Definitions of
Variance}\label{mathematical-definitions-of-variance}}

Let \(A\) and \(B\) be types where \(A\) is a subtype of \(B\)
(\(A \subseteq B\)). Let \(F\) be a generic container type constructor
(e.g., \passthrough{\lstinline!List[T]!},
\passthrough{\lstinline!Iterable[T]!}, or
\passthrough{\lstinline!Callable[[T], R]!}).

\begin{itemize}
\tightlist
\item
  \textbf{Covariance}: The container subtype relationship preserves the
  parameter type relationship: \[F(A) \subseteq F(B)\] \emph{Example}:
  \passthrough{\lstinline!Iterable[T]!} is covariant. A list of dogs can
  be safely read as a list of animals.
\item
  \textbf{Contravariance}: The container subtype relationship reverses
  the parameter type relationship: \[F(B) \subseteq F(A)\]
  \emph{Example}: \passthrough{\lstinline!Callable[[T], None]!} is
  contravariant. A function that accepts any animal can safely be used
  where a function accepting a dog is expected.
\item
  \textbf{Invariance}: There is no relationship between container types:
  \[F(A) \not\subseteq F(B) \quad \text{and} \quad F(B) \not\subseteq F(A)\]
  \emph{Example}: \passthrough{\lstinline!List[T]!} is invariant because
  it allows both read (covariant) and write (contravariant) operations.
  Allowing a list of dogs to be treated as a list of animals could let
  someone insert a cat into the list, violating type safety.
\end{itemize}

\begin{center}\rule{0.5\linewidth}{0.5pt}\end{center}
